Throughout this chapter, let $C$ be a smooth, projective, geometrically connected
curve over $k$. Fix a modulus

\[
    \mdls = \sum_i n_i P_i
\]

on $C$, and set $U = C \setminus \mdls$.

For each single-point submodulus $\ndls = n_i P_i$ of $\mdls$, set
$d_\ndls = \deg(\ndls)$. The effective divisor $\ndls$ determines a
$k$-rational closed point

\[
    Z_\ndls \hookrightarrow C^{(d_\ndls)}
        = \operatorname{Sym}^{d_\ndls}_k(C).
\]

Let

\[
    \pi_\ndls \colon
    X_\ndls := \operatorname{Bl}_{Z_\ndls}\bigl(C^{(d_\ndls)}\bigr)
    \longrightarrow C^{(d_\ndls)}
\]

be the blowup at this point. We denote its exceptional divisor by

\[
    E_\ndls := Z_\ndls \times_{C^{(d_\ndls)}} X_\ndls
\]

By
\Cref{prop:relative_symmetric_power_smooth,prop:exceptional_divisor_smooth_point},
$E_\ndls\cong\mathbb P_k^{d_\ndls-1}$; in particular, it is irreducible and
has codimension $1$ in $X_\ndls$. We denote its unique generic point by
$\eta_\ndls$, so that $\overline{\{\eta_\ndls\}}=E_\ndls$. Thus we have the
Cartesian diagram

\begin{equation}\label{equation:blowup_diagram_ndls}
\begin{tikzcd}
E_\ndls = \overline{\{\eta_\ndls\}}
    \arrow[r, hook] \arrow[d]
  & X_\ndls \arrow[d, "\pi_\ndls"] \\
Z_\ndls \arrow[r, hook]
  & C^{(d_\ndls)}.
\end{tikzcd}
\end{equation}

The purpose of this chapter is to prove
\Cref{theorem:SymmetricPowerOfSheavesIsTamelyRamifiedReduction}. Let $\cF$ be
the rank-one $\Lambda$-local system appearing there and put
$\cP=\operatorname{Fr}(\cF)$. The theorem becomes the following torsor
statement.\footnote{Here one uses
$\operatorname{Fr}(\cF^{[d_\ndls]})\cong\cP^{[d_\ndls]}$ and the fact that a
local system and its frame torsor have the same inertia character, as
recorded in a footnote in Chapter~1.}

\begin{theorem}\label{theorem:torsors_after_blowup_is_tame}
Let $G$ be a finite abelian group, and let $\cP \to U$ be a $G$-torsor whose
ramification is bounded by $\mdls$. For every single-point submodulus
$\ndls = n_iP_i$ of $\mdls$, the symmetric-power torsor
$\cP^{[d_\ndls]}$ is tamely ramified at $\eta_\ndls$.
\end{theorem}

After recalling the necessary local ramification theory, we make two
reductions. First, the general finite abelian case is reduced to cyclic groups.
Second, formal--local invariance reduces the geometric calculation on $C$ to a
calculation on $\mathbb{G}_m \subset \mathbb{P}^1$. These reductions motivate the
Kummer, Artin--Schreier, and Artin--Schreier--Witt theories used in the cyclic
calculations that follow.
